\documentclass[11pt]{amsart}
\usepackage[margin=1in]{geometry}
\usepackage{amssymb}
\usepackage[no-math]{fontspec}
\setmonofont{DejaVu Sans Mono}[Scale=0.72]
\usepackage{longtable}
\usepackage{url}
\usepackage[hidelinks]{hyperref}
\newtheorem{theorem}{Theorem}[section]
\newtheorem{proposition}[theorem]{Proposition}
\theoremstyle{definition}
\newtheorem{definition}[theorem]{Definition}
\theoremstyle{remark}
\newtheorem{remark}[theorem]{Remark}
\title[NoetherLefschetz]{Noether--Lefschetz Vanishing for Hypersurfaces in $\mathbb P^5$ as an Explicit Hypothesis: a Machine-Checked Lean~4 Interface}
\author{Benjamin Stanley Frohman}
\email{frohmanbenjamin@gmail.com}
\urladdr{https://github.com/BenFrohman/NoetherLefschetz}
\date{November 4, 2026}
\sloppy
\begin{document}
\begin{abstract}
\noindent\textbf{DOI:} \texttt{10.5281/zenodo.23942455}\par\smallskip

We record in Lean~4, with no user axioms, the propositional form of Noether--Lefschetz vanishing: a very general smooth hypersurface of degree $d\ge3$ in $\mathbb P^5$ has no rational Hodge class in $H^4$ beyond multiples of $h^2$. Hypersurfaces are abstract records with a degree, a ``very general'' predicate and an ``extra class'' predicate. Conditional on the vanishing sentence, we prove that an extra class on such a hypersurface of degree at least $3$ gives a contradiction, and that there is no extra class when the degree is at least $6$. A second encoding over an abstract moduli setting proves that the $d\ge3$ sentence implies its $d\ge6$ case. The predicate $\mathrm{Hdg}^2=\mathbb Q h^2$ in that encoding is a placeholder defined as \texttt{True}. All results are conditional. Beauville's monodromy theorem and Deligne's invariance argument are not formalized.
\end{abstract}
\maketitle

\section{Introduction}
For a very general smooth hypersurface $X\subset\mathbb P^5$ of degree $d\ge3$, the space of rational Hodge classes $H^4(X,\mathbb Q)\cap H^{2,2}(X)$ is spanned by $h^2$. The standard proof combines Beauville's theorem that monodromy is Zariski dense in the orthogonal group of vanishing cohomology~\cite{beauville} with Deligne's theorem that a Hodge class on a very general fibre is monodromy invariant~\cite{deligne}. For $d=2$ the statement fails, because smooth quadric fourfolds contain planes. Neither argument is formalized here. The repository previously asserted the vanishing (for $d\ge6$) and the moduli-level statement (for $d\ge3$) as Lean axioms, and one file did not compile. In this version every axiom is removed. The vanishing is a hypothesis, and all results are implications from it.

\section{Encoding by records (\texttt{NoetherLefschetz.Basic})}
\begin{definition}[\texttt{Hypersurface}, \texttt{ExtraClass}, \texttt{VanishingOfExtraClasses}]
A \emph{hypersurface record} $X$ has a degree $\deg X\in\mathbb N$ and propositions \texttt{isVeryGeneral} and \texttt{hasExtraClass}. \texttt{ExtraClass X} is \texttt{hasExtraClass}. The \emph{vanishing sentence} for $X$ is: if $3\le\deg X$ and $X$ is very general, then $X$ has no extra class. The proposition \texttt{noether\_lefschetz X} is defined to be this sentence. \texttt{HodgeLocusPosAlgebraic d} is defined to be $6\le d$.
\end{definition}
\begin{theorem}[Conditional; \texttt{no\_extra\_on\_nl\_locus}]
Let $X$ satisfy the vanishing sentence. If $3\le\deg X$, $X$ is very general and $X$ has an extra class, then \texttt{False}.
\end{theorem}
\begin{theorem}[Conditional; \texttt{noether\_lefschetz\_of\_vanishing}]
Let $X$ satisfy the vanishing sentence. If $6\le\deg X$ and $X$ is very general, then $X$ has no extra class.
\end{theorem}

\section{Encoding over a moduli setting (\texttt{NoetherLefschetz.Spec})}
\begin{definition}[\texttt{SpecSetting}, \texttt{hdg2\_eq\_Qh2}, \texttt{extra\_rational\_classes\_vanish}]
A \emph{spec setting} $S$ consists of types $M_d$ for $d\in\mathbb N$ (moduli of smooth degree-$d$ hypersurfaces) and a predicate \texttt{IsVeryGeneral} on properties of points of $M_d$. The predicate \texttt{hdg2\_eq\_Qh2} is a placeholder \emph{defined as} \texttt{True}. The sentence \texttt{extra\_rational\_classes\_vanish d} is: $3\le d$ implies $\texttt{IsVeryGeneral}(X\mapsto\texttt{hdg2\_eq\_Qh2}\,d\,X)$.
\end{definition}
\begin{proposition}[\texttt{extra\_rational\_classes\_vanish\_unfolds}]
For every $S$ and $d$, \texttt{extra\_rational\_classes\_vanish d} equals its defining formula.
\end{proposition}
\begin{theorem}[Conditional; \texttt{d\_ge\_6\_subcase}]
Assume that for all $d\ge3$, $\texttt{IsVeryGeneral}(X\mapsto\texttt{hdg2\_eq\_Qh2}\,d\,X)$. Then the same holds for every $d\ge6$.
\end{theorem}
\begin{remark}
Because \texttt{hdg2\_eq\_Qh2} is \texttt{True}, the second encoding carries no Hodge-theoretic content. It records only how the threshold $d\ge6$ follows from $d\ge3$.
\end{remark}

\section{Correspondence between claims and Lean declarations}\label{sec:table}
Every mathematical claim of this paper is listed below with the Lean declaration that proves it. ``Conditional'' means the hypothesis named in the statement is an explicit argument of the Lean theorem and is \emph{not} proved in the repository. No claim in this paper goes beyond these statements.

{\footnotesize
\begin{longtable}{p{0.30\linewidth}p{0.40\linewidth}p{0.22\linewidth}}
\hline
Claim & Lean declaration & Status \\
\hline\endhead
Vanishing, $\deg\ge3$, very general, extra class $\Rightarrow$ False & \path{NoetherLefschetz.no_extra_on_nl_locus} & Conditional on \texttt{VanishingOfExtraClasses X} \\
Vanishing $\Rightarrow$ no extra class for $\deg\ge6$ & \path{NoetherLefschetz.noether_lefschetz_of_vanishing} & Conditional on \texttt{VanishingOfExtraClasses X} \\
Sentence unfolds to its definition & \path{NoetherLefschetz.Spec.extra_rational_classes_vanish_unfolds} & Proved (by definition) \\
$d\ge3$ sentence $\Rightarrow$ $d\ge6$ case & \path{NoetherLefschetz.Spec.d_ge_6_subcase} & Conditional on the $d\ge3$ hypothesis \\
\hline
\end{longtable}}

\section{Scope and limitations}
Hodge theory, monodromy, moduli of hypersurfaces and the theorems of Beauville and Deligne are not formalized. Hypersurface records and spec settings are arbitrary data, and \texttt{hdg2\_eq\_Qh2} is \texttt{True}. The Noether--Lefschetz theorem is not proved; the results are implications from it. The bound $6$ is not claimed to be sharp. Statement changes in this version: the axioms \texttt{ExtraClass} and \texttt{noether\_lefschetz} in \texttt{Basic} are removed. \texttt{Basic} did not compile and has been repaired: \texttt{hasExtraClass} is a record field, and \texttt{ExtraClass} and \texttt{noether\_lefschetz} are definitions. The old axiom \texttt{noether\_lefschetz} ($d\ge6$) is now the conditional theorem \texttt{noether\_lefschetz\_of\_vanishing}. In \texttt{Spec}, the axioms \texttt{SmoothHypersurfaceModuli} and \texttt{IsVeryGeneral} are now fields of \texttt{SpecSetting}, and the axiom \texttt{extra\_rational\_classes\_vanish\_holds} is now a hypothesis of \texttt{d\_ge\_6\_subcase}.


\appendix
\section{Lean source}
The complete Lean source of the certified development (copyright headers omitted) is reproduced verbatim.
\subsection*{\texttt{NoetherLefschetz/Basic.lean}}
\begin{verbatim}
/-!
# Noether–Lefschetz as named propositions

Not the Hodge conjecture. Literature, not a proof in this file.

Geometric range: a very general smooth hypersurface of degree `d ≥ 3` in `ℙ⁵`
has no extra rational Hodge class. Degree `2` fails (quadric fourfolds contain planes).
No `axiom`. No `sorry`. No dummy `P ∧ ¬P`.
`VanishingOfExtraClasses` is a sentence. A term of that sentence would be
Beauville (monodromy dense in the orthogonal group of vanishing cohomology)
plus Deligne (a rational Hodge class on a very general fibre is
monodromy-invariant). Those arguments are not in this file; every result is
conditional on that sentence, supplied as a hypothesis.
-/

namespace NoetherLefschetz

/-- A named hypersurface in ℙ⁵, as a record.
`hasExtraClass` is the geometric predicate “Hdg² larger than ℚ h²”.
It is a field, not `degree = 0 ∧ degree ≠ 0`. -/
structure Hypersurface where
  degree : Nat
  isVeryGeneral : Prop
  hasExtraClass : Prop

/-- Extra means not a rational multiple of h². -/
def ExtraClass (X : Hypersurface) : Prop :=
  X.hasExtraClass

/-- Vanishing sentence (monodromy + Deligne).
Very general smooth hypersurface in ℙ⁵ of degree ≥ 3. -/
def VanishingOfExtraClasses (X : Hypersurface) : Prop :=
  3 ≤ X.degree → X.isVeryGeneral → ¬ ExtraClass X

/-- BKU sentence, as a name only: the inequality `6 ≤ d`. -/
def HodgeLocusPosAlgebraic (d : Nat) : Prop :=
  6 ≤ d

/-- Former axiom name, now a proposition. -/
def noether_lefschetz (X : Hypersurface) : Prop :=
  VanishingOfExtraClasses X

/-- Implication from the named vanishing sentence. Does not inhabit that sentence. -/
theorem no_extra_on_nl_locus (X : Hypersurface)
    (hV : VanishingOfExtraClasses X)
    (hd : 3 ≤ X.degree) (hvg : X.isVeryGeneral)
    (hExtra : ExtraClass X) : False :=
  hV hd hvg hExtra

/-- The former axiom `noether_lefschetz` (degree ≥ 6), now conditional on the
vanishing sentence. -/
theorem noether_lefschetz_of_vanishing (X : Hypersurface)
    (hV : VanishingOfExtraClasses X)
    (hd : 6 ≤ X.degree) (hvg : X.isVeryGeneral) :
    ¬ ExtraClass X :=
  hV (Nat.le_trans (by decide) hd) hvg

end NoetherLefschetz
\end{verbatim}
\subsection*{\texttt{NoetherLefschetz/Spec.lean}}
\begin{verbatim}
/-
Statement encoding. No Mathlib imports.
The former axioms `SmoothHypersurfaceModuli`, `IsVeryGeneral` are fields of
`SpecSetting`; the former axiom `extra_rational_classes_vanish_holds` is a hypothesis.
Does not prove Beauville, Deligne, or Noether-Lefschetz.
Note: `hdg2_eq_Qh2` is a placeholder defined as `True`.
-/

namespace NoetherLefschetz
namespace Spec

/-- Abstract moduli of smooth degree-`d` hypersurfaces in P^5 and a
"very general" predicate. Supplied by the user. -/
structure SpecSetting where
  SmoothHypersurfaceModuli : Nat → Type
  IsVeryGeneral : {d : Nat} → (SmoothHypersurfaceModuli d → Prop) → Prop

variable (S : SpecSetting)

/-- H^4(X,Q) ∩ H^{2,2}(X) = Q h^2. Placeholder: defined as `True`. -/
def hdg2_eq_Qh2 (d : Nat) (_X : S.SmoothHypersurfaceModuli d) : Prop :=
  True

/-- The sentence. Proof is not supplied. -/
def extra_rational_classes_vanish (d : Nat) : Prop :=
  3 ≤ d → S.IsVeryGeneral fun X => hdg2_eq_Qh2 S d X

theorem extra_rational_classes_vanish_unfolds (d : Nat) :
    extra_rational_classes_vanish S d =
      (3 ≤ d → S.IsVeryGeneral fun X => hdg2_eq_Qh2 S d X) :=
  rfl

/-- d ≥ 6 is a subcase of the d ≥ 3 sentence (conditional on that sentence). -/
theorem d_ge_6_subcase
    (extra_rational_classes_vanish_holds : ∀ d, 3 ≤ d → S.IsVeryGeneral fun X => hdg2_eq_Qh2 S d X)
    (d : Nat) (hd : 6 ≤ d) :
    S.IsVeryGeneral fun X => hdg2_eq_Qh2 S d X :=
  extra_rational_classes_vanish_holds d (Nat.le_trans (by decide) hd)

end Spec
end NoetherLefschetz
\end{verbatim}

\section{Verification certificate}\label{app:verification-certificate}

The formal development was machine-checked as follows.

{\raggedright
\begin{itemize}
\item Repository:\\ {\footnotesize\url{https://github.com/BenFrohman/NoetherLefschetz}}
\item Verified commit:\\ {\footnotesize\texttt{ec7289969d31f50bf84395a515875466a656422d}}
\item Lean toolchain: \texttt{leanprover/lean4:v4.34.0}
\item Mathlib: none (the development depends only on Lean core)
\item Continuous integration run (build, axiom audit, \texttt{sorry}/\texttt{axiom} scan; passed):\\ {\footnotesize\url{https://github.com/BenFrohman/NoetherLefschetz/actions/runs/41794912612}}
\item SHA-256 of the GitHub source tarball of the verified commit, downloaded twice with identical result:\\ {\footnotesize\url{https://github.com/BenFrohman/NoetherLefschetz/archive/ec7289969d31f50bf84395a515875466a656422d.tar.gz}}\\
{\footnotesize\texttt{43d91047d05a339cb5ce78982e0f84b4d44582dadeea0a82efccd5db78dc565f}}
\end{itemize}\par}

The tarball hash refers to the verified commit rather than to a release tag, since a document cannot contain the hash of an archive that includes itself. The Lean sources are unchanged between that commit and the release tag.

\textbf{The development contains no \texttt{sorry} and no user-declared axioms.} Every one of the 4 theorems and lemmas listed below depends only on Lean's standard axioms \texttt{propext}, \texttt{Classical.choice} and \texttt{Quot.sound} (or on none). The table is the verbatim output of \texttt{lake env lean AxiomAudit.lean} (\texttt{\#print axioms}) on the verified commit.

{\footnotesize
\begin{longtable}{p{0.55\linewidth}p{0.38\linewidth}}
\hline
Lean declaration & \texttt{\#print axioms} \\
\hline\endhead
\path{NoetherLefschetz.no_extra_on_nl_locus} & \texttt{(no axioms)} \\
\path{NoetherLefschetz.noether_lefschetz_of_vanishing} & \texttt{(no axioms)} \\
\path{NoetherLefschetz.Spec.d_ge_6_subcase} & \texttt{(no axioms)} \\
\path{NoetherLefschetz.Spec.extra_rational_classes_vanish_unfolds} & \texttt{(no axioms)} \\
\hline
\end{longtable}}

\begin{thebibliography}{9}
\bibitem{beauville} A.~Beauville, \emph{Le groupe de monodromie des familles universelles d'hypersurfaces et d'intersections compl\`etes}, Lecture Notes in Math. 1194, Springer (1986), 8--18.
\bibitem{deligne} P.~Deligne, \emph{Th\'eorie de Hodge II}, Publ. Math. IH\'ES 40 (1971), 5--57.
\bibitem{lean} L.~de Moura and S.~Ullrich, \emph{The Lean 4 theorem prover and programming language}, CADE-28, LNCS 12699 (2021), 625--635.
\end{thebibliography}

\end{document}
